% Optimisation and Dykstra's algorithm
\addsymbol{$d$}{Optimisation state dimension}
\addsymbol{$n$}{Total number of half-spaces or particles}
\addsymbol{$\mathcal{H}$}{Polyhedral feasible set formed by the intersection of half-spaces}
\addsymbol{$A$}{Constraint matrix}
\addsymbol{$a_i$}{Unit normal vector defining the boundary of the $i$-th half-space}
\addsymbol{$b_i$}{Boundary offset parameter of the $i$-th half-space}
\addsymbol{$w$}{Spatial coordinate decision variable}
\addsymbol{$w^\circ$}{Initial unconstrained point prior to projection}
\addsymbol{$w^\star$}{True Euclidean projection onto the feasible set}
\addsymbol{$m$}{Inner Dykstra projection cycle index}
\addsymbol{$t$}{Outer projected gradient descent (PGD) learning loop index}
\addsymbol{$w^{(t)}$}{Spatial weights at outer iteration $t$}
\addsymbol{$w^{(m)}$}{Primary spatial coordinate at inner Dykstra iteration $m$}
\addsymbol{$w^{(t, m)}$}{Double-indexed weights representing inner state $m$ at outer step $t$}
\addsymbol{$e^{(m)}$}{Auxiliary memory variable evaluated at inner iteration $m$}
\addsymbol{$k^{(m)}$}{Active half-space constraint index at inner iteration $m$}
\addsymbol{$M^{(t)}$}{Dynamic inner Dykstra cycle iteration budget ceiling}
\addsymbol{$M_{\mathrm{base}}$}{Base projection budget for the inner cycle}
\addsymbol{$p$}{Dynamic iteration budget scaling exponent}
\addsymbol{$\epsilon_m$}{Monotonicity tolerance vector}
\addsymbol{$N_{\mathrm{stall}}$}{Closed-form stall duration}
\addsymbol{$m_{\mathrm{stall}}$}{Iteration index denoting the onset of the stalling phenomenon}
\addsymbol{$N_{\mathrm{skip}}$}{Fast-forward algorithmic skip duration}
\addsymbol{$\mathcal{B}^{(t)}$}{Random subset of indices for stochastic mini-batching at iteration $t$}
\addsymbol{$B$}{Mini-batch cardinality}
\addsymbol{$g^{(t)}$}{Stochastic gradient evaluated over the selected mini-batch}
\addsymbol{$\eta^{(t)}$}{Optimisation learning rate at iteration $t$}
\addsymbol{$\eta^\circ$}{Initial base learning rate}
\addsymbol{$\gamma$}{Learning rate decay parameter}
\addsymbol{$\mathcal{M}^{(t)}$}{Boolean active mask vector governing iterative hard thresholding}
\addsymbol{$\epsilon_p$}{Threshold parameter for pruning microscopic weights in IHT}
\addsymbol{$K$}{Iterative hard thresholding application interval}
\addsymbol{$T$}{Total outer PGD iteration limit}

% Optimal transport and hypothesis space
\addsymbol{$D$}{Physical state dimension of the underlying dynamical system}
\addsymbol{$J_k + 1$}{Total dimension of the polynomial coefficient space}
\addsymbol{$P$}{Maximum individual or total polynomial degree}
\addsymbol{$\alpha^{(j)}$}{Multi-index governing the $j$-th multivariate basis function}
\addsymbol{$\alpha_m^{(j)}$}{The $m$-th spatial component of the multi-index $\alpha^{(j)}$}
\addsymbol{$x$}{Unmapped target particles (sheared representation)}
\addsymbol{$z$}{Mapped particles targeted to a standard normal distribution}
\addsymbol{$\tilde{z}$}{True multivariate standard normal reference samples}
\addsymbol{$\tilde{S}$}{Artificial non-linear shear mapping function}
\addsymbol{$\mathcal{S}_k$}{Function hypothesis space for scalar map components}
\addsymbol{$\ell_2$}{Standard Euclidean distance norm}

% Orbital mechanics and dynamics
\addsymbol{$a$}{Semi-major axis of an orbital trajectory}
\addsymbol{$e$}{Eccentricity of the orbit}
\addsymbol{$i$}{Orbital inclination relative to the reference plane}
\addsymbol{$\Omega$}{Right ascension of the ascending node}
\addsymbol{$\omega$}{Argument of periapsis}
\addsymbol{$\nu$}{True anomaly defining specific angular position}
\addsymbol{$h, k$}{Components of the non-singular eccentricity vector}
\addsymbol{$p, q$}{Components of the non-singular ascending node vector}
\addsymbol{$\lambda$}{Mean longitude}
\addsymbol{$\mathbb{R}^5 \times \mathbb{S}^1$}{Topological space parameterising the equinoctial orbital elements}